\documentclass[11pt,a4paper]{article}
\usepackage[T1]{fontenc}
\usepackage[utf8]{inputenc}
\usepackage{lmodern}
\usepackage[a4paper,margin=23mm]{geometry}
\usepackage{amsmath,amssymb,bm,booktabs,longtable,array}
\usepackage{microtype}
\usepackage{xcolor}
\usepackage[hidelinks]{hyperref}
\usepackage{url}
\hypersetup{pdftitle={Model electric-field gradients in anilite: analytic derivation and numerical conventions}}
\setlength{\parindent}{1.2em}
\setlength{\parskip}{3pt}
\setlength{\emergencystretch}{2em}
\renewcommand{\theequation}{S\arabic{equation}}
\newcommand{\dd}{\mathrm d}
\newcommand{\erfc}{\operatorname{erfc}}
\newcommand{\erf}{\operatorname{erf}}
\newcommand{\Tr}{\operatorname{Tr}}
\newcommand{\diag}{\operatorname{diag}}
\newcommand{\Id}{\mathbf 1}
\newcommand{\G}{\bm G}
\newcommand{\V}{\bm V}
\newcommand{\rvec}{\bm r}
\newcommand{\kvec}{\bm k}
\newcommand{\nvec}{\bm n}
\newcommand{\nuq}{\nu_Q}
\newcommand{\nunqr}{\nu_{\mathrm{NQR}}}
\newcommand{\nuz}{\nu_{\mathrm Z}}
\newcommand{\etapc}{\eta_{\mathrm{pc}}}
\begin{document}
\lefthyphenmin=3
\righthyphenmin=3
\hyphenation{NQR NMR EFG EFGs}
\begin{center}
{\LARGE\bfseries Model electric-field gradients in anilite:\par}
\vspace{3pt}
{\Large analytic derivation, numerical conventions,\par
and the low-frequency Cu3 hypothesis\par}
\vspace{9pt}
{\small Supporting calculation for the copper powder-NMR manuscript\par
27 November 2026\par}
\end{center}

\noindent\textbf{Scope.} This document derives the electrostatic lattice proxy
used for anilite Cu$_7$S$_4$ and the steps relating an EFG to spin-$3/2$ NQR and
powder NMR. It retains the distinction between (i) a converged mathematical
lattice sum, (ii) a model of the complete nuclear EFG, and (iii) a crystallographic
assignment of an observed signal. The calculations establish the first of these
for explicitly chosen point charges. They provide hypotheses for the other two.
Symbolic intermediate expressions and checks are saved by
\path{theory/symbolic_algebra.py}; all real-space and reciprocal-space numerical
terms are retained separately under \path{theory/results/efg/}.

\setcounter{tocdepth}{1}
\tableofcontents
\newpage

\section{Definitions and the structural input}

Let a periodic unit cell contain charges $q_j=z_j e$ at Cartesian positions
$\rvec_j=\bm A\bm f_j$, where $\bm f_j$ are fractional coordinates. For the
reference orthorhombic anilite structure,
\begin{equation}
 \bm A=\diag(7.89,7.84,11.01)\ \text{\AA},\qquad
 \Omega=\det\bm A=681.052176\ \text{\AA}^3.
\end{equation}
The value computed from the listed lattice lengths is used rather than the
rounded CIF volume. The coordinate source is the transcription of AMCSD~0009376
preserved in \path{research_literature/anilite_amcsd0009376_transcribed.cif},
based on Koto and Morimoto's structure \cite{Koto1970}. It is a reference
near-room-temperature structure, not a refinement of the measured sample at 5~K.

\begin{center}
\begin{tabular}{lrrrl}
\toprule
Atom & $f_a$ & $f_b$ & $f_c$ & Unit-cell multiplicity \\
\midrule
Cu1 & 0.48320 & 0.75000 & 0.09420 & 4 \\
Cu2 & 0.42580 & 0.25000 & 0.33350 & 4 \\
Cu3 & 0.47900 & 0.75000 & 0.35450 & 4 \\
Cu4 & 0.23710 & 0.43030 & 0.15530 & 8 \\
Cu5 & 0.17700 & 0.47670 & 0.41860 & 8 \\
S1 & 0.24600 & 0.25000 & 0.98730 & 4 \\
S2 & 0.78500 & 0.25000 & 0.98410 & 4 \\
S3 & 0.99880 & 0.99800 & 0.26110 & 8 \\
\bottomrule
\end{tabular}
\end{center}

Each representative coordinate is expanded by the eight symmetry operations
stored in the CIF, reduced modulo the unit cell, and deduplicated. This gives
28 Cu and 16 S atoms, or four formula units. Cu1 and Cu3 are tetrahedrally
coordinated; Cu2, Cu4 and Cu5 have three nearest sulfur neighbours. The count of
eight tetrahedrally and twenty triangularly coordinated Cu atoms per unit cell
agrees with the published crystallochemical comparison \cite{Evans1981}.

Two types of charge scenario are evaluated. The baseline uses
\begin{equation}
 z_{\mathrm{Cu}}=\frac87,\qquad z_{\mathrm S}=-2,\qquad
 28\frac87-16(2)=0.
\label{eq:neutral}
\end{equation}
Five further scenarios start with $z_{\mathrm{Cu}}=1$ and place the remaining
four positive charge units on one complete Cu orbit of multiplicity $m_s$:
\begin{equation}
 z_s=1+\frac4{m_s},\quad z_{\text{other Cu}}=1,
 \quad z_{\mathrm S}=-2.
\end{equation}
All six cells are neutral and retain the reference space-group symmetry. These
are sensitivity scenarios, not measured oxidation-state assignments.

\section{From the Coulomb potential to the EFG kernel}

\subsection{Potential Hessian and sign convention}

For a point charge at a displacement $\rvec$ from the observation point, the
electrostatic potential is
\begin{equation}
 \phi(\rvec)=\frac{ze}{4\pi\epsilon_0 r},\qquad
 r=(r_a r_a)^{1/2}.
\end{equation}
Repeated Cartesian indices are summed in this section. We use the conventional
quadrupolar EFG sign: $V_{ab}$ is the second derivative of the electrostatic
\emph{potential}. Since $\bm E=-\nabla\phi$, the derivative
$\partial E_a/\partial r_b$ has the opposite sign. A consistent convention is
essential when retaining signed eigenvalues, even though the zero-field NQR
frequency below involves their magnitude.

First differentiate $r$ and $1/r$:
\begin{align}
 \frac{\partial r}{\partial r_a}&=\frac{r_a}{r},\\
 \frac{\partial (r^{-1})}{\partial r_a}&=-\frac{r_a}{r^3},\\
 \frac{\partial^2(r^{-1})}{\partial r_a\partial r_b}
 &=-\frac{\delta_{ab}}{r^3}+\frac{3r_a r_b}{r^5}
 =\frac{3r_a r_b-r^2\delta_{ab}}{r^5}.
\label{eq:kernel}
\end{align}
Hence the numerical lattice tensor used in the code is
\begin{equation}
 G_{ab}(s)=\sum_{j,\bm R}'z_j
 \frac{3d_{j\bm R,a}d_{j\bm R,b}
 -d_{j\bm R}^2\delta_{ab}}{d_{j\bm R}^5},\qquad
 \bm d_{j\bm R}=\rvec_j+\bm R-\rvec_s.
\label{eq:direct}
\end{equation}
The prime removes the observation atom in the home cell, while its nonzero
periodic images remain. The kernel is even in $\bm d$, so defining the
displacement with the opposite sign gives the same Hessian.

For each nonzero displacement,
\begin{equation}
 \sum_a\frac{3d_a^2-d^2}{d^5}=0.
\end{equation}
Therefore the lattice tensor is symmetric and traceless. For a general continuous
charge distribution the quadrupolar interaction selects the traceless part of
the potential Hessian. Here the observation point is away from all retained
point-charge singularities, and the lattice Hessian is already traceless.

\subsection{Units and what an SI conversion means}

The scripts use $z_j=q_j/e$ as numerical charge values and distances in
\AA. Thus the numerical entries of $G$ are in \AA$^{-3}$ for charges measured
in units of $e$. Some archived files describe this convention as
``$e$/\AA$^3$, with the Coulomb prefactor omitted.'' There is exactly one factor
of the elementary charge in the SI conversion:
\begin{equation}
 V^{\mathrm{pc}}_{ab},[\mathrm{V\,m^{-2}}]
 =\underbrace{\frac{e}{4\pi\epsilon_0}
 (10^{10})^3}_{\simeq\,1.43996\times10^{21}\ \mathrm{V\,m^{-2}}}
 G_{ab}\,[\text{numerical units}].
\label{eq:units}
\end{equation}
The constants follow the 2022 CODATA convention \cite{CODATA}. Equation
\eqref{eq:units} is a unit conversion for the chosen point-charge model. It does
not restore the omitted electronic charge density or make the model a prediction
of the complete nuclear EFG.

\section{Periodic Ewald sum, including the self correction}

\subsection{Potential splitting}

The Ewald identity is
\begin{equation}
 \frac1r=\frac{\erfc(\alpha r)}r+\frac{\erf(\alpha r)}r,
 \qquad \alpha>0.
\end{equation}
It partitions the periodic Coulomb problem into a short-range real-space sum
and a smooth reciprocal-space sum \cite{Ewald1921,DLPOLY}. The reciprocal vectors
satisfy $\kvec=2\pi\bm A^{-T}\bm m$, with integer $\bm m$. For a neutral cell,
the nonzero reciprocal modes of the smooth potential, without the common SI
prefactor, are
\begin{equation}
 \psi^{\mathrm{rec}}(\rvec)=\frac{4\pi}{\Omega}
 \sum_{\kvec\ne0}\frac{e^{-k^2/(4\alpha^2)}}{k^2}
 \sum_j z_j\cos[\kvec\cdot(\rvec-\rvec_j)].
\label{eq:recpotential}
\end{equation}
Charge neutrality removes the problematic uniform charge mode. An arbitrary
constant potential does not affect its Hessian. The present calculation uses
the periodic electrostatic convention; it does not introduce a spatially varying
macroscopic external field gradient.

\subsection{Real-space derivatives, with all radial coefficients}

For any radial function $f(r)$,
\begin{align}
 \partial_a f&=f'(r)\frac{r_a}{r},\\
 \partial_a\partial_b f
 &=\left(f''-\frac{f'}r\right)\frac{r_a r_b}{r^2}
 +\frac{f'}r\delta_{ab}.
\label{eq:radial}
\end{align}
Set $f(r)=\erfc(\alpha r)/r$ and $E(r)=e^{-\alpha^2r^2}$. Differentiation gives
\begin{align}
 f'(r)&=-\frac{\erfc(\alpha r)}{r^2}
 -\frac{2\alpha}{\sqrt\pi}\frac{E(r)}r,\\
 f''(r)&=\frac{2\erfc(\alpha r)}{r^3}
 +\frac{4\alpha}{\sqrt\pi}\frac{E(r)}{r^2}
 +\frac{4\alpha^3}{\sqrt\pi}E(r).
\end{align}
Substituting these two expressions into Eq.~\eqref{eq:radial} gives
\begin{align}
 A(r)&=\frac{3\erfc(\alpha r)}{r^5}
 +\frac{6\alpha E(r)}{\sqrt\pi r^4}
 +\frac{4\alpha^3E(r)}{\sqrt\pi r^2},\\
 B(r)&=-\frac{\erfc(\alpha r)}{r^3}
 -\frac{2\alpha E(r)}{\sqrt\pi r^2},\\
 G^{\mathrm{real}}_{ab}(s)
 &=\sum_{j,\bm R}'z_j
 \left[A(d)d_a d_b+B(d)\delta_{ab}\right].
\label{eq:real}
\end{align}
The trace of one screened Hessian is
\begin{equation}
 A(r)r^2+3B(r)=\frac{4\alpha^3}{\sqrt\pi}E(r).
\label{eq:realtrace}
\end{equation}
Thus the real-space Ewald subtotal is \emph{not} separately traceless. Discarding
its trace before adding the other contributions would alter the required
cancellation.

\subsection{Reciprocal-space Hessian}

Twice differentiating the cosine in Eq.~\eqref{eq:recpotential} gives
\begin{equation}
 \partial_a\partial_b\cos[\kvec\cdot(\rvec-\rvec_j)]
 =-k_a k_b\cos[\kvec\cdot(\rvec-\rvec_j)].
\end{equation}
Consequently,
\begin{equation}
 G^{\mathrm{rec}}_{ab}(s)=-\frac{4\pi}{\Omega}
 \sum_{\kvec\ne0}\frac{e^{-k^2/(4\alpha^2)}}{k^2}k_a k_b
 \sum_jz_j\cos[\kvec\cdot(\rvec_s-\rvec_j)].
\label{eq:rec}
\end{equation}
The sum as implemented includes both $\kvec$ and $-\kvec$. No additional factor
of two is applied. A half-space implementation would need the corresponding
factor explicitly.

\subsection{Why the self correction has a positive Hessian}

The observation charge is removed from the singular real-space sum, but its
smooth Gaussian contribution is still contained in Eq.~\eqref{eq:rec}. Near
that charge,
\begin{equation}
 \frac{\erf(\alpha r)}r
 =\frac{2\alpha}{\sqrt\pi}
 -\frac{2\alpha^3}{3\sqrt\pi}r^2
 +\frac{\alpha^5}{5\sqrt\pi}r^4+O(r^6).
\end{equation}
The constant has zero Hessian and
$\partial_a\partial_b r^2=2\delta_{ab}$. Hence the smooth self term already
present in the reciprocal sum has Hessian
\begin{equation}
 G^{\mathrm{smooth\ self}}_{ab}
 =-\frac{4\alpha^3z_s}{3\sqrt\pi}\delta_{ab}.
\end{equation}
Removing it therefore \emph{adds} the opposite tensor:
\begin{equation}
 \boxed{\quad
 \G(s)=\G^{\mathrm{real}}(s)+\G^{\mathrm{rec}}(s)
 +\frac{4\alpha^3z_s}{3\sqrt\pi}\Id .\quad}
\label{eq:ewald}
\end{equation}
This correction follows by differentiating with respect to the observation
coordinate while holding source positions fixed. Differentiating a constant
self-energy correction would not derive the local field Hessian correctly.

In exact arithmetic Eq.~\eqref{eq:ewald} is traceless. The code records the trace
\emph{before} projecting out a residual $\Tr(\G)\Id/3$. That final projection
only removes roundoff, after the real, reciprocal and self terms have been
combined; it cannot replace the self correction.

\subsection{Convergence is distinct from model validity}

The archived calculation uses $r_c=6/\alpha$, $k_c=12\alpha$, with
$\alpha=0.2,0.3,0.4$~\AA$^{-1}$. At these cutoffs the exponential factors reach
$e^{-36}$. Agreement across $\alpha$ tests the cancellation between two different
partitions of the same lattice sum. Direct spherical sums at radii
15, 20, 30, 40, 60, 80 and 100~\AA\ provide a separate convergence comparison.
The primary periodic tensors change by less than approximately
$3\times10^{-15}$ numerical tensor units between the three splitting parameters.
These are numerical checks for a fixed model; they are not an uncertainty on
the physical nuclear EFG.

\section{Principal axes, asymmetry, and symmetry restrictions}

Diagonalize the real symmetric tensor:
\begin{equation}
 \G\bm u_i=\lambda_i\bm u_i,\qquad
 \bm U^T\G\bm U=\diag(G_{XX},G_{YY},G_{ZZ}).
\end{equation}
Choose $|G_{ZZ}|\geq|G_{YY}|\geq|G_{XX}|$. Tracelessness then permits the
convention
\begin{equation}
 \etapc=\frac{G_{XX}-G_{YY}}{G_{ZZ}},\qquad0\leq\etapc\leq1.
\end{equation}
The reverse reconstruction follows from the sum and difference:
\begin{equation}
 G_{XX}=\frac{\etapc-1}{2}G_{ZZ},\qquad
 G_{YY}=-\frac{\etapc+1}{2}G_{ZZ}.
\end{equation}
Eigenvector signs are arbitrary; eigenvector order is not. Near a degenerate
pair the numerical orientation of its individual eigenvectors is not uniquely
defined, even when the tensor itself is well determined. If $G_{ZZ}=0$, the
ordering and trace imply $\G=0$, and $\etapc$ is undefined rather than zero.

For a $4c$ site of this $Pnma$ setting, reflection perpendicular to $b$ is
$\bm M=\diag(1,-1,1)$. Invariance gives
\begin{equation}
 \bm M\G\bm M^T-\G=
 \begin{pmatrix}0&-2G_{ab}&0\\-2G_{ab}&0&-2G_{bc}\\0&-2G_{bc}&0\end{pmatrix}
 =0.
\end{equation}
Thus $G_{ab}=G_{bc}=0$ and $b$ is one principal direction. The $ac$ block may
contain a nonzero off-diagonal entry, and neither the mirror nor the general
$8d$ site symmetry forces $\eta=0$.

For four equal charges at the vertices proportional to
$(1,1,1)$, $(1,-1,-1)$, $(-1,1,-1)$ and $(-1,-1,1)$,
\begin{equation}
 \sum_j \left(3\bm d_j\bm d_j^T-d_j^2\Id\right)=0.
\end{equation}
An ideal tetrahedron therefore gives a vanishing tensor. A distorted tetrahedron
can give a small residual tensor with substantial asymmetry. Small $|G_{ZZ}|$
and small $\etapc$ are separate statements.

\section{Worked Cu3 calculation and the five-site comparison}

For the baseline charges, the periodic Cu3 tensor in crystallographic axes is
\begin{equation}
 \G_{\mathrm{Cu3}}=10^{-3}
 \begin{pmatrix}
 11.66175955&0&31.28442026\\
 0&-14.47092006&0\\
 31.28442026&0&2.80916052
 \end{pmatrix}.
\end{equation}
Entries suppressed to zero are below the displayed precision and respect the
mirror constraint. Diagonalization and ordering by magnitude give
\begin{equation}
 (G_{XX},G_{YY},G_{ZZ})
 =(-0.01447092006,-0.02436053779,0.03883145786).
\end{equation}
The trace vanishes within rounding, and the resulting asymmetry is
\begin{equation}
 \etapc(\mathrm{Cu3})=
 \frac{-0.01447092006+0.02436053779}{0.03883145786}
 =0.25468057.
\end{equation}
Multiplying this tensor by Eq.~\eqref{eq:units} changes its units but leaves its
eigenvectors, eigenvalue ratios and asymmetry unchanged.

\begin{center}
\begin{tabular}{lrrr}
\toprule
Site & Signed $G_{ZZ}$ & $\etapc$ & $|G_{ZZ}|/|G_{ZZ}^{\mathrm{Cu2}}|$\\
\midrule
Cu1 & $-0.30264647$ & 0.55719893 & 0.68774730 \\
Cu2 & $ 0.44005476$ & 0.05914175 & 1.00000000 \\
Cu3 & $ 0.03883146$ & 0.25468057 & 0.08824233 \\
Cu4 & $ 0.37291139$ & 0.16291879 & 0.84742041 \\
Cu5 & $ 0.41868874$ & 0.58492897 & 0.95144691 \\
\bottomrule
\end{tabular}
\end{center}

Cu3 has the smallest tensor magnitude in all six neutral scenarios. This
property extends throughout their restricted convex charge family: the tensor
depends linearly on charge, so its spectral norm obeys
\begin{equation}
 \left\|\sum_iw_i\G_i\right\|_2
 \leq\sum_iw_i\|\G_i\|_2
 \leq\max_i\|\G_i\|_2,\qquad w_i\geq0,\quad\sum_iw_i=1.
\end{equation}
For Cu3 the upper bound from the five charge-redistribution vertices is 0.08088.
For each other Cu site a projection along its baseline principal direction has
one sign at all five vertices, with magnitude at least 0.28488, 0.40150, 0.34181
or 0.38664 for Cu1, Cu2, Cu4 or Cu5, respectively. A signed projection bounds the
spectral norm from below, so none of these sites can become smaller than Cu3
within that charge simplex. This is a property of the specified fixed-coordinate
model family, not a theorem about arbitrary electronic charge distributions.

\section{From the EFG to the \texorpdfstring{spin-$3/2$}{spin-3/2} NQR frequency}

\subsection{Quadrupolar Hamiltonian and conventions}

Use dimensionless angular-momentum operators, with
$I_Z|m\rangle=m|m\rangle$ and
\begin{equation}
 I_\pm|m\rangle=
 \sqrt{I(I+1)-m(m\pm1)}\,|m\pm1\rangle.
\end{equation}
For the physical nuclear EFG, the quadrupole Hamiltonian is
\begin{align}
 \frac{\mathcal H_Q}{h}
 &=\frac{\nuq}{6}
 \left[3I_Z^2-I(I+1)+\eta(I_X^2-I_Y^2)\right],\\
 \nuq&=\frac{3eQV_{ZZ}}{2I(2I-1)h}.
\label{eq:nuqdef}
\end{align}
For $I=3/2$, $\nuq=eQV_{ZZ}/(2h)$. The signed $\nuq$ used here is a coupling
constant; it should not be confused with a paper's notation in which the symbol
$\nu_Q$ denotes the experimentally observed NQR frequency itself. The standard
Cu NQR Hamiltonian and this identifiability issue are discussed in
Ref.~\cite{Gainov2009}.

Since $I_X^2-I_Y^2=(I_+^2+I_-^2)/2$, in the ordered basis
$\{|3/2\rangle,|1/2\rangle,|-1/2\rangle,|-3/2\rangle\}$,
\begin{equation}
 \frac{\mathcal H_Q}{h}=\frac{\nuq}{2}
 \begin{pmatrix}
 1&0&\eta/\sqrt3&0\\
 0&-1&0&\eta/\sqrt3\\
 \eta/\sqrt3&0&-1&0\\
 0&\eta/\sqrt3&0&1
 \end{pmatrix}.
\end{equation}
The two $2\times2$ blocks have the same characteristic polynomial:
\begin{equation}
 \det\left[
 \frac{\nuq}{2}
 \begin{pmatrix}1&\eta/\sqrt3\\\eta/\sqrt3&-1\end{pmatrix}
 -\lambda\Id\right]
 =\lambda^2-\frac{\nuq^2}{4}\left(1+\frac{\eta^2}{3}\right).
\end{equation}
There are two doubly degenerate energy levels,
\begin{equation}
 E_\pm/h=\pm\frac{|\nuq|}{2}\sqrt{1+\frac{\eta^2}{3}},
\end{equation}
and their separation gives
\begin{equation}
 \boxed{\nunqr=|\nuq|\sqrt{1+\frac{\eta^2}{3}}.}
\label{eq:nqr}
\end{equation}

\subsection{What frequency measurements identify}

One spin-$3/2$ NQR frequency constrains a combination of $V_{ZZ}$ and $\eta$.
For any chosen trial asymmetry,
\begin{equation}
 |\nuq|=\frac{\nunqr}{\sqrt{1+\eta^2/3}},\qquad
 \frac{\sqrt3}{2}\nunqr\leq|\nuq|\leq\nunqr.
\end{equation}
For two Cu isotopes in the same environment, neglecting an isotope-induced
change of the electronic environment,
\begin{equation}
 \frac{{}^{65}\nunqr}{{}^{63}\nunqr}
 =\left|\frac{{}^{65}Q}{{}^{63}Q}\right|.
\end{equation}
The second isotope is useful for pairing lines but does not supply a new
independent equation separating $V_{ZZ}$ from $\eta$ in zero field. In applied
field its different gyromagnetic ratio provides a further check on the combined
Zeeman--quadrupole profile.

For the thesis frequencies 11.97 and 16.96~MHz, the bound gives
\begin{equation}
 \frac{|V_{ZZ,s4}|}{|V_{ZZ,s1}|}
 \geq\frac{(\sqrt3/2)16.96}{11.97}=1.22705\ldots.
\end{equation}
Thus the four high-frequency groups cannot all be explained by equal
$|V_{ZZ}|$ with only different asymmetries.

\section{Finite-field levels and the unoriented-powder calculation}

For the isotropic magnetic-shift model used in the archived fit,
\begin{equation}
 \frac{\mathcal H(\nvec)}h=-\nuz\,\bm I\cdot\nvec
 +\frac{\nuq}{6}
 [3I_Z^2-I(I+1)+\eta(I_X^2-I_Y^2)],
\label{eq:fullH}
\end{equation}
where $\nvec$ is the field direction in the EFG principal frame and
$\nuz=(\gamma/2\pi)B_0(1+K)$ for a scalar shift $K$. Each trial $\eta$ is
accompanied by the $|\nuq|$ required by Eq.~\eqref{eq:nqr}. The physical sign is
not established by the present analysis; the archived fits use the positive
coupling magnitude.

Diagonalizing the $4\times4$ matrix gives eigenvectors $|u\rangle$ and
eigenvalues $\epsilon_u=E_u/h$. The transition frequencies are
$\nu_{uv}=\epsilon_v-\epsilon_u>0$. For a transverse RF direction $\bm e$,
the matrix-element weight is $|\langle u|\bm I\cdot\bm e|v\rangle|^2$.
Averaging its azimuth about $\nvec$ uses
\begin{equation}
 \langle e_a e_b\rangle_\perp
 =\frac12(\delta_{ab}-n_a n_b),
\end{equation}
which gives
\begin{equation}
 W_{uv}(\nvec)=\frac12\left[
 \sum_a|\langle u|I_a|v\rangle|^2
 -|\langle u|\bm I\cdot\nvec|v\rangle|^2\right].
\end{equation}
The corresponding orientation-averaged profile is
\begin{equation}
 P(\nu)=\frac1{4\pi}\int_0^{2\pi}\!\dd\varphi
 \int_{-1}^{1}\!\dd u\,
 \sum_{u_1<v_1}W_{u_1v_1}(u,\varphi)
 \delta[\nu-\nu_{u_1v_1}(u,\varphi)],
\label{eq:powder}
\end{equation}
where the integration variable $u=\cos\theta$ is distinct from an eigenstate
label. No preferred orientation is imposed. For Eq.~\eqref{eq:fullH}, symmetry
allows the integral to be reduced to one octant. Equal-spaced midpoint grids in
$\cos\theta$ and $\varphi$ then have equal weights after normalization.

Numerical frequency deposition is followed by convolution with a Gaussian of
standard deviation $\sigma$. The measured spectrum is modelled as a sum over
four spectral groups and two isotopes, with group amplitudes and an empirical
underlying contribution. The 20 and 500~ms repetition-period spectra share
their quadrupolar parameters but may have different amplitudes. These weights
describe a frequency-profile model; actual spin-echo intensities additionally
depend on populations, pulse excitation, RF distribution, relaxation and probe
response. They are not direct crystallographic occupancies.

The powder shape depends on $\eta$ through both angular variations in
Eq.~\eqref{eq:fullH} and the required change of $\nuq$ at fixed $\nunqr$.
Nevertheless, overlapping profiles and unmodelled response factors can leave
multiple acceptable minima. In particular, the archived fit retains low- and
high-$\eta_1$ branches; a point-charge model does not by itself remove that
experimental identifiability problem.

\section{Why a small-EFG site need not create distant wings}

\subsection{Explicit high-field derivation for an axial EFG}

For an illustrative analytic limit, set $\eta=0$ and place the quantization
axis along $B_0$. The EFG axis is then represented by
$I_Z^{\mathrm{EFG}}=I_z\cos\theta+I_x\sin\theta$. The unperturbed energies are
\begin{equation}
 E_m^{(0)}/h=-\nuz m.
\end{equation}
Writing $c=\cos\theta$, first-order perturbation theory gives
\begin{equation}
 E_m^{(1)}/h=\frac{\nuq}{12}(3c^2-1)
 [3m^2-I(I+1)].
\end{equation}
The $m=1/2$ and $m=-1/2$ shifts are equal, so the central transition has no
first-order quadrupole shift. Second-order perturbation theory gives
\begin{equation}
 E_m^{(2)}/h=\sum_{m'\ne m}
 \frac{|(\mathcal H_Q/h)_{mm'}|^2}{\nuz(m'-m)}.
\end{equation}
For $m=1/2$, the two nonzero couplings have squared magnitudes
$3\nuq^2(1-c^2)c^2/4$ and $3\nuq^2(1-c^2)^2/16$ and denominators
$\nuz$ and $-2\nuz$, respectively. Therefore
\begin{align}
 E_{1/2}^{(2)}/h
 &=\frac{3\nuq^2(1-c^2)c^2}{4\nuz}
 -\frac{3\nuq^2(1-c^2)^2}{32\nuz},\\
 E_{-1/2}^{(2)}/h&=-E_{1/2}^{(2)}/h,\\
 \delta\nu_{\mathrm{CT}}^{(2)}
 &=\frac{3\nuq^2}{16\nuz}(1-c^2)(1-9c^2).
\label{eq:ct}
\end{align}
Maximizing and minimizing the quadratic in $c^2\in[0,1]$ gives
\begin{equation}
 -\frac{\nuq^2}{3\nuz}\leq\delta\nu_{\mathrm{CT}}^{(2)}
 \leq\frac{3\nuq^2}{16\nuz}.
\end{equation}
The lower bound is reached at $c^2=5/9$ and the upper bound at $c^2=0$.
The first-order satellite frequencies lie at
\begin{equation}
 \nu_{\mathrm{sat}}=\nuz\pm\frac{\nuq}{2}(3c^2-1),
\end{equation}
so their combined outer bounds are $\nuz\pm|\nuq|$ in this approximation.

\subsection{A numerical scale example, not a Cu3 frequency measurement}

For $\nuq=1.5$~MHz, $\nuz=76$~MHz and $\eta=0$, the central-transition bounds
from Eq.~\eqref{eq:ct} are $-9.868$ and $+5.551$~kHz; the satellite bounds are
approximately 74.5--77.5~MHz. Such a signal can overlap the already occupied
$^{63}$Cu frequency region. Absence of additional resolved outer wings therefore
constrains a proposed fifth contribution but does not exclude every low-EFG
fifth site. The frequency 1.5~MHz is an illustrative input, not an assignment
deduced from this high-field expansion. Exact diagonalization is used in the
archived numerical profile calculations.

\section{Empirical scaling of the lattice proxy and the low-frequency band}

\subsection{An explicit conditional test}

To test whether the \emph{relative} lattice tensors are compatible with the
frequency scale, define
\begin{equation}
 P_s=|G_{ZZ,s}|\sqrt{1+\eta_{\mathrm{pc},s}^2/3}.
\end{equation}
Assume provisionally that one positive common scale $a$ maps $P_s$ to NQR
frequencies. This amounts to assuming a common multiplicative relation between
the nuclear EFGs and these lattice proxies across sites. It is an additional
physical hypothesis; the electrostatic calculation does not establish it.

For each of the $4!=24$ assignments $\pi$ of Cu1, Cu2, Cu4 and Cu5 to the four
high-frequency groups $\nu_i$, minimize
\begin{align}
 L_\pi(a)&=\sum_{i=1}^4[aP_{\pi(i)}-\nu_i]^2,\\
 \frac{\dd L_\pi}{\dd a}
 &=2\sum_iP_{\pi(i)}[aP_{\pi(i)}-\nu_i]=0,\\
 a_\pi&=\frac{\sum_iP_{\pi(i)}\nu_i}
 {\sum_iP_{\pi(i)}^2},\\
 \widehat\nu_{\mathrm{Cu3},\pi}&=a_\pi P_{\mathrm{Cu3}}.
\label{eq:calibration}
\end{align}
All 24 assignments are retained for each of the six charge scenarios, including
their high-frequency residuals. The low-frequency band is not used to fit $a$.
The calibration uses the thesis high-frequency set
$(11.97,15.27,16.18,16.96)$~MHz. A frequency-only optimum is not a simultaneous
fit of the experimental asymmetries.

For the baseline point charges, the best frequency-only mapping is
$(s1,s2,s3,s4)\mapsto(\mathrm{Cu1,Cu4,Cu2,Cu5})$, with
$a=38.29166696$~MHz per numerical proxy unit. Its explicit residuals are
\begin{center}
\begin{tabular}{lrrr}
\toprule
Group & Input frequency (MHz) & Scaled proxy (MHz) & Residual (MHz)\\
\midrule
s1 & 11.970000 & 12.173743 & $+0.203743$\\
s2 & 15.270000 & 14.342428 & $-0.927572$\\
s3 & 16.180000 & 16.860251 & $+0.680251$\\
s4 & 16.960000 & 16.921832 & $-0.038168$\\
\bottomrule
\end{tabular}
\end{center}
The high-frequency RMSE is 0.584401~MHz. The scaled Cu3 value is
1.502909~MHz. The mapped baseline asymmetries are approximately
$(0.557,0.163,0.059,0.585)$; they do not reproduce the fitted powder-NMR
asymmetry pattern. Hence even the best frequency-only mapping does not validate
a unique assignment.

\begin{center}
\begin{tabular}{lrr}
\toprule
Charge scenario & Best-scaled Cu3 (MHz) & High-frequency RMSE (MHz)\\
\midrule
Uniform Cu charge $8/7$ & 1.50291 & 0.5844\\
Excess charge on Cu1 & 3.51559 & 0.5922\\
Excess charge on Cu2 & 2.18006 & 0.3713\\
Excess charge on Cu3 & 1.97518 & 0.5229\\
Excess charge on Cu4 & 1.40027 & 0.8775\\
Excess charge on Cu5 & 1.45307 & 0.7256\\
\bottomrule
\end{tabular}
\end{center}

These selected-scenario values span approximately 1.40--3.52~MHz. They are
conditional scaled estimates, not an independently predicted absolute NQR
frequency and not a confidence interval. The complete 144-assignment table is
saved, so alternative assignments and poor frequency matches are not hidden.

\subsection{What the historical low-frequency observation establishes}

The author's thesis reports a signal in the 1.5--2.2~MHz region at 4.2~K for a
mixed sample of nominal overall composition Cu$_{1.1}$S (p.~22, Fig.~8), and
discusses tetrahedral environments and unresolved phase contributions later in
the text \cite{Thesis}. This is evidence of an observed low-frequency response,
not a resolved assignment to a particular anilite site. The experimenter's
current account of limited tuning and resonant-circuit coverage provides
additional experimental context; it is not substituted for an archived
frequency-resolved dataset.

A concrete competing contribution is already known: covellite CuS has an
experimentally observed $^{63}$Cu NQR signal at 1.87~MHz and 4.2~K
\cite{Gainov2009}. Accordingly, a Cu3 contribution from anilite is a possible
component of the historical mixed-sample response. The existence of a band near
2~MHz alone cannot separate it from CuS. The lattice-proxy result motivates a
test of that component; it does not assign the entire low-frequency band to
anilite.

\section{Physical limitations and a path to a quantitative assignment}

The complete nuclear EFG contains the actual electronic charge distribution.
Schematically one may distinguish a lattice contribution and an on-site
electronic contribution. Semiempirical treatments sometimes write \cite{Gainov2009}
\begin{equation}
 \V^{\mathrm{nuc}}\simeq(1-\gamma_\infty)\V^{\mathrm{lattice}}
 +(1-R)\V^{\mathrm{valence}},
\end{equation}
with antishielding and shielding parameters. This is only a schematic
decomposition here; no values of $\gamma_\infty$, $R$, orbital populations or
site-independent enhancement factor have been assumed or fitted. Covalency,
metallic screening and site-dependent electronic effects can change the
eigenvalues, axes and asymmetries. A scalar rescaling changes tensor magnitudes
but cannot change $\eta$, and therefore cannot repair an asymmetry mismatch.

Thermal displacements introduce another distinction. For a relative
displacement $\bm u_j-\bm u_s$, the thermal average of a nonlinear kernel
$K_{ab}(\bm d)$ generally obeys
\begin{equation}
 \langle K_{ab}(\bm d+\bm u)\rangle
 =K_{ab}(\bm d)+\frac12\partial_\mu\partial_\nu K_{ab}(\bm d)
 \langle u_\mu u_\nu\rangle+\cdots,
\end{equation}
rather than simply the kernel at the average coordinate. Relative-displacement
covariances require both individual atomic displacements and their correlations.
The CIF displacement parameters alone do not determine a complete thermal EFG
or $\eta(T)$. Likewise,
\begin{equation}
 \frac{\dd\ln\nunqr}{\dd T}
 =\frac{\dd\ln|V_{ZZ}|}{\dd T}
 +\frac{\eta}{3+\eta^2}\frac{\dd\eta}{\dd T}
\end{equation}
does not separate those two contributions using only one NQR frequency.

The present model's robust result is the comparatively small Cu3 lattice tensor
within its stated charge family. A quantitative experimental assignment would
require a reliable low-temperature structure, an electronic-structure EFG
calculation, and a joint match of all available frequencies, asymmetries,
intensity constraints and phase contributions. The four high-frequency spectral
labels s1--s4 are retained separately from Cu1--Cu5 throughout the numerical
archive.

\section{Reproducibility and saved intermediates}

The repository preserves the earlier calculation unchanged and adds an
independent, more detailed evaluation. The principal entry points are:
\begin{itemize}
\item \path{theory/symbolic_algebra.py}: Coulomb derivatives, screened radial
derivatives, reciprocal sign, smooth-self Taylor series, mirror and tetrahedron
identities, explicit spin matrices, zero-field eigenvalues, all second-order
state shifts, and the analytic common-scale least-squares solution.
\item In \path{theory/results/}, the files \path{symbolic_algebra.json} and
\path{symbolic_algebra.txt} retain every saved symbolic step in plain
and LaTeX forms; \path{symbolic_checks.csv} records each identity residual.
\item \path{theory/results/efg/}: expanded cell, per-term real and reciprocal
contributions, self terms, subtotal tensors, raw traces, eigensystems,
convergence results, and all empirical frequency-scaling assignments.
\item \path{new_evidence/structure_model/geometry_efg_proxy.py}: original
periodic point-charge calculation and its retained results.
\end{itemize}
The detailed numerical file names and commands are listed in the accompanying
Russian readme. All algebraic checks compare independently derived expressions,
and the numerical checks compare summation methods or exact Hamiltonian results;
they do not merely assert that copied output matches itself.

\begin{thebibliography}{9}
\bibitem{Koto1970}
K. Koto and N. Morimoto, The crystal structure of anilite,
\emph{Acta Crystallographica B} \textbf{26}, 915--924 (1970).
\href{https://doi.org/10.1107/S0567740870003370}{doi:10.1107/S0567740870003370}.
The coordinate transcription follows AMCSD~0009376 as exposed at
\url{https://www.mindat.org/show.php?clearcache=1&id=236&ld=}.

\bibitem{Evans1981}
H. T. Evans, Jr., Copper coordination in low chalcocite and djurleite and other
copper-rich sulfides, \emph{American Mineralogist} \textbf{66}, 807--818 (1981).

\bibitem{Ewald1921}
P. P. Ewald, Die Berechnung optischer und elektrostatischer Gitterpotentiale,
\emph{Annalen der Physik} \textbf{369}, 253--287 (1921).
\href{https://doi.org/10.1002/andp.19213690304}{doi:10.1002/andp.19213690304}.

\bibitem{DLPOLY}
DL\_Software Digital Guide, Distance cutoffs and Ewald sums;
DL\_POLY manual, Choosing Ewald Sum Variables.
\url{https://dl-sdg.github.io/RESOURCES/FORCE_FIELD/ff11.html};
\url{https://ccp5.gitlab.io/dl-poly/UserManual/construction_and_execution/ewald_precision.html}.
Accessed 27 November 2026.

\bibitem{Gainov2009}
R. R. Gainov \emph{et al.}, Phase transition and anomalous electronic behavior
in the layered superconductor CuS probed by NQR,
\emph{Physical Review B} \textbf{79}, 075115 (2009).
\href{https://doi.org/10.1103/PhysRevB.79.075115}{doi:10.1103/PhysRevB.79.075115};
\url{https://arxiv.org/abs/0809.0560}.

\bibitem{Thesis}
I. A. Evlampiev, \emph{Nuclear quadrupole resonance of copper in geerite},
master's thesis, Kazan State University (2009), in Russian.
Author-provided file \path{diploma_magister_geerite_2009_5.doc};
low-frequency discussion on pp.~22 and 33. The historical phase label is
preserved in the source title.

\bibitem{CODATA}
NIST, The 2022 CODATA Recommended Values of the Fundamental Physical Constants,
Web Version 9.0 (2024). \url{https://physics.nist.gov/constants}.
Elementary charge $e=1.602176634\times10^{-19}$~C (exact);
$\epsilon_0=8.8541878188(14)\times10^{-12}$~F\,m$^{-1}$.
\end{thebibliography}
\end{document}
